\nonumchap{Abstract}

\draftnote{replace with a 200--300 word abstract once the results chapters are finalised.}

This thesis investigates the detection of subwavelength displacement and
material change in ground-penetrating radar (GPR) time-lapse surveys. Classical
migration imaging is shown to be amplitude-limited: point scatterers separated,
or displaced, by a fraction of a wavelength cannot be resolved from the migrated
amplitude image alone. A 2D Fourier phase-plane shift-estimation method is
developed instead, which exploits the Fourier shift theorem applied to the
cross-spectrum of baseline and monitor migrated images. A weighted least-squares
fit of the cross-spectrum phase recovers sub-millimetre lateral and vertical
displacements, and decouples purely geometric movement from a change in
sub-wavelength material properties (e.g.\ a fracture filling with fluid). The
method is validated on synthetic gprMax models across three migration
algorithms (Kirchhoff, Gazdag phase-shift, and time-reversal back-propagation),
across displacement scales from $2\lambda$ down to $\tfrac{1}{32}\lambda$, under
additive noise, and on a more realistic distributed fluid-front scenario.
